%! Adding \& Subtracting Rational Expressions — Unit 4, Lesson 4.3 — Homework
\documentclass[10pt]{article}
\usepackage{algebra2-article}
\usepackage{algebra2-boxes}
\usepackage{grading}
\gradingassignment{alg2-u04-l03-hw}
\gradingcourse{algebra_2}
\usepackage{algebra2-graded}
\newcommand{\TallMath}[1]{$\displaystyle #1\rule[-1.4em]{0pt}{3.2em}$}

\begin{document}

\pageheader{Unit 4, Lesson 4.3}{Homework}

\vspace{0.10in}

\begin{notesbox}{Practice}
Every answer needs \emph{two} parts: the simplest form \textbf{and} the restrictions. Factor every
denominator before you build an LCD, and put the second numerator in \textbf{parentheses} every single
time you subtract. Write each final answer in its box. Only what is inside a box is graded.

\begin{enumerate}[leftmargin=1.9em, itemsep=11pt]
  \item \textbf{Like denominators.} Combine, simplify, restrict.
    \begin{enumerate}[label=(\alph*), leftmargin=2.1em, itemsep=5pt]
      \item $\dfrac{3x}{x+7}+\dfrac{21}{x+7}=\;$\inblank{\blank{2.2cm}}{\gradedblank[id=q1a, type=short, points=2, answer={3}, check=numeric]{2.2cm-2pt}{18pt}} \quad $x\neq$ \inblank{\blank{1.4cm}}{\gradedblank[id=q1b, type=short, points=1, answer={-7}, check=numeric]{1.4cm-2pt}{18pt}}
      \item $\dfrac{7x-1}{x^2-1}-\dfrac{3x+3}{x^2-1}=\;$\inblank{\blank{2.2cm}}{\gradedblank[id=q1c, type=short, points=2, answer={4/(x+1)}, display={$\frac{4}{x+1}$}]{2.2cm-2pt}{18pt}} \quad $x\neq$ \inblank{\blank{1.8cm}}{\gradedblank[id=q1d, type=short, points=1, answer={1, -1}, check=set, display={$1,\ -1$}]{1.8cm-2pt}{18pt}}
      \item $\dfrac{x^2}{x+5}-\dfrac{25}{x+5}=\;$\inblank{\blank{2.2cm}}{\gradedblank[id=q1e, type=short, points=2, answer={x-5}]{2.2cm-2pt}{18pt}} \quad $x\neq$ \inblank{\blank{1.4cm}}{\gradedblank[id=q1f, type=short, points=1, answer={-5}, check=numeric]{1.4cm-2pt}{18pt}}
    \end{enumerate}
    \vspace{2pt}
    Your answer to (a) is a number and your answer to (c) is a polynomial. Neither has a denominator
    left. So what does each restriction do to the graph? \rubricbox{\gradedblank[id=q1g, type=rubric, points=2, display={it punches a hole: (a) is the horizontal line $y=3$ missing $(-7,3)$, and (c) is the line $y=x-5$ missing $(-5,-10)$}]{\linewidth}{10pt}}

  \item \textbf{Monomial LCDs.}
    \begin{enumerate}[label=(\alph*), leftmargin=2.1em, itemsep=5pt]
      \item $\dfrac{3}{4x}+\dfrac{5}{6x}=\;$\inblank{\blank{2.2cm}}{\gradedblank[id=q2a, type=short, points=2, answer={19/(12x)}, display={$\frac{19}{12x}$}]{2.2cm-2pt}{18pt}} \quad LCD used: \inblank{\blank{1.6cm}}{\gradedblank[id=q2b, type=short, points=1, answer={12x}]{1.6cm-2pt}{18pt}}
      \item $\dfrac{7}{x^2}-\dfrac{2}{5x}=\;$\inblank{\blank{2.2cm}}{\gradedblank[id=q2c, type=short, points=2, answer={(35-2x)/(5x^2)}, display={$\frac{35-2x}{5x^2}$}]{2.2cm-2pt}{18pt}} \quad LCD used: \inblank{\blank{1.6cm}}{\gradedblank[id=q2d, type=short, points=1, answer={5x^2}]{1.6cm-2pt}{18pt}}
      \item $\dfrac{4}{3a}+\dfrac{1}{a^2b}=\;$\inblank{\blank{2.2cm}}{\gradedblank[id=q2e, type=short, points=2, answer={(4ab+3)/(3a^2b)}, display={$\frac{4ab+3}{3a^2b}$}]{2.2cm-2pt}{18pt}} \quad restrictions: \inblank{\blank{2.2cm}}{\gradedblank[id=q2f, type=short, points=1, answer={a!=0, b!=0}, check=set, display={$a\neq0,\ b\neq0$}]{2.2cm-2pt}{18pt}}
    \end{enumerate}

  \item \textbf{Binomial and quadratic LCDs.} Factor first. One of these has \emph{opposite}
        binomials for denominators, and one of them cancels at the end.
    \begin{enumerate}[label=(\alph*), leftmargin=2.1em, itemsep=7pt]
      \item $\dfrac{2}{x+1}+\dfrac{3}{x-2}=\;$\inblank{\blank{2.6cm}}{\gradedblank[id=q3a, type=short, points=2, answer={(5x-1)/((x+1)(x-2))}, display={$\frac{5x-1}{(x+1)(x-2)}$}]{2.6cm-2pt}{18pt}} \quad $x\neq$ \inblank{\blank{1.8cm}}{\gradedblank[id=q3b, type=short, points=1, answer={-1, 2}, check=set, display={$-1,\ 2$}]{1.8cm-2pt}{18pt}}
      \item $\dfrac{x}{x-3}-\dfrac{2}{x+3}=\;$\inblank{\blank{2.6cm}}{\gradedblank[id=q3c, type=short, points=2, answer={(x^2+x+6)/((x-3)(x+3))}, display={$\frac{x^2+x+6}{(x-3)(x+3)}$}]{2.6cm-2pt}{18pt}} \quad $x\neq$ \inblank{\blank{1.8cm}}{\gradedblank[id=q3d, type=short, points=1, answer={3, -3}, check=set, display={$3,\ -3$}]{1.8cm-2pt}{18pt}}
      \item $\dfrac{4}{x-6}-\dfrac{48}{x^2-36}=\;$\inblank{\blank{2.6cm}}{\gradedblank[id=q3e, type=short, points=2, answer={4/(x+6)}, display={$\frac{4}{x+6}$}]{2.6cm-2pt}{18pt}} \quad $x\neq$ \inblank{\blank{1.8cm}}{\gradedblank[id=q3f, type=short, points=1, answer={6, -6}, check=set, display={$6,\ -6$}]{1.8cm-2pt}{18pt}}
      \item $\dfrac{6}{x-1}+\dfrac{3}{1-x}=\;$\inblank{\blank{2.6cm}}{\gradedblank[id=q3g, type=short, points=2, answer={3/(x-1)}, display={$\frac{3}{x-1}$}]{2.6cm-2pt}{18pt}} \quad $x\neq$ \inblank{\blank{1.4cm}}{\gradedblank[id=q3h, type=short, points=1, answer={1}, check=numeric]{1.4cm-2pt}{18pt}}
      \item $\dfrac{5}{x^2+4x+3}+\dfrac{2}{x^2-9}=\;$\inblank{\blank{2.6cm}}{\gradedblank[id=q3i, type=short, points=2, answer={(7x-13)/((x+1)(x+3)(x-3))}, display={$\frac{7x-13}{(x+1)(x+3)(x-3)}$}]{2.6cm-2pt}{18pt}} \quad $x\neq$ \inblank{\blank{2.2cm}}{\gradedblank[id=q3j, type=short, points=1, answer={-1, -3, 3}, check=set, display={$-1,\ -3,\ 3$}]{2.2cm-2pt}{18pt}}
    \end{enumerate}

  \item \textbf{Complex fractions.} Simplify each as a quotient of two simple fractions, then state
        \emph{all} restrictions --- the inner denominators \emph{and} the main fraction bar.
    \begin{enumerate}[label=(\alph*), leftmargin=2.1em, itemsep=7pt]
      \item \TallMath{\frac{\;\dfrac{1}{x}+\dfrac{1}{6}\;}{\;\dfrac{1}{x}-\dfrac{1}{6}\;}}
            $=\;$\inblank{\blank{2.2cm}}{\gradedblank[id=q4a, type=short, points=2, answer={(6+x)/(6-x)}, display={$\frac{6+x}{6-x}$}]{2.2cm-2pt}{18pt}} \quad $x\neq$ \inblank{\blank{1.8cm}}{\gradedblank[id=q4b, type=short, points=1, answer={0, 6}, check=set, display={$0,\ 6$}]{1.8cm-2pt}{18pt}}
      \item \TallMath{\frac{\;\dfrac{2}{x}-\dfrac{2}{5}\;}{\;x-5\;}}
            $=\;$\inblank{\blank{2.2cm}}{\gradedblank[id=q4c, type=short, points=2, answer={-2/(5x)}, display={$-\frac{2}{5x}$}]{2.2cm-2pt}{18pt}} \quad $x\neq$ \inblank{\blank{1.8cm}}{\gradedblank[id=q4d, type=short, points=1, answer={0, 5}, check=set, display={$0,\ 5$}]{1.8cm-2pt}{18pt}}
      \item \TallMath{\frac{\;\dfrac{x}{x+2}\;}{\;\dfrac{x^2}{x^2-4}\;}}
            $=\;$\inblank{\blank{2.2cm}}{\gradedblank[id=q4e, type=short, points=2, answer={(x-2)/x}, display={$\frac{x-2}{x}$}]{2.2cm-2pt}{18pt}} \quad $x\neq$ \inblank{\blank{2.2cm}}{\gradedblank[id=q4f, type=short, points=1, answer={0, 2, -2}, check=set, display={$0,\ 2,\ -2$}]{2.2cm-2pt}{18pt}}
    \end{enumerate}

  \item \textbf{A round trip.} A cyclist rides $30$ miles out to a lake at $x$ mi/h and rides the same
        $30$ miles home at $x+5$ mi/h.
    \begin{enumerate}[label=(\alph*), leftmargin=2.1em, itemsep=5pt]
      \item Total time $T(x)=$ \inblank{\blank{4.4cm}}{\gradedblank[id=q5a, type=short, points=2, answer={30(2x+5)/(x(x+5))}, display={$30(2x+5)/\bigl(x(x+5)\bigr)$}]{4.4cm-2pt}{10pt}} hours (one fraction, simplified).
      \item Check your formula: if $x=10$, part (a) gives \inblank{\blank{1.2cm}}{\gradedblank[id=q5b, type=short, points=1, answer={5}, check=numeric]{1.2cm-2pt}{10pt}} hours, and adding the two legs
            by hand gives \inblank{\blank{1.2cm}}{\gradedblank[id=q5c, type=short, points=1, answer={5}, check=numeric]{1.2cm-2pt}{10pt}} hours.
      \item The algebra excludes two values of $x$. Which two? \inblank{\blank{1.8cm}}{\gradedblank[id=q5d, type=short, points=1, answer={0, -5}, check=set, display={$0$ and $-5$}]{1.8cm-2pt}{10pt}} \; Which values of $x$
            make sense for a bike ride? \inblank{\blank{1.8cm}}{\gradedblank[id=q5e, type=short, points=1, answer={x>0}]{1.8cm-2pt}{10pt}}
      \item \textbf{Average speed} for the whole trip is total distance divided by total time --- a
            complex fraction. Simplify $\dfrac{60}{T(x)}=$ \inblank{\blank{3.0cm}}{\gradedblank[id=q5f, type=short, points=2, answer={2x(x+5)/(2x+5)}, display={$\frac{2x(x+5)}{2x+5}$}]{3.0cm-2pt}{18pt}} mi/h. \\*[3pt]
            At $x=10$ that gives \inblank{\blank{1.4cm}}{\gradedblank[id=q5g, type=short, points=1, answer={12}, check=numeric]{1.4cm-2pt}{10pt}} mi/h. Is that the average of $10$ and $15$?
            \inblank{\blank{1.0cm}}{\gradedblank[id=q5h, type=short, points=1, answer={no}, check=text]{1.0cm-2pt}{10pt}} \; Explain in one sentence. \rubricbox{\gradedblank[id=q5i, type=rubric, points=2, display={you spend more \emph{time} at the slower speed, so it counts for more}]{\linewidth}{10pt}}
    \end{enumerate}

  \item \textbf{Justify.} Are $\dfrac{1}{x-2}+\dfrac{1}{x+2}$ and $\dfrac{2x}{x^2-4}$ the same
        expression? Then answer the same question for $\dfrac{3}{x-3}-\dfrac{18}{x^2-9}$ and
        $\dfrac{3}{x+3}$. Answer both in terms of \emph{inputs}: which inputs does each form accept,
        and where do they agree?
        \gradedlines{3}{\gradedblank[id=q6, type=rubric, points=4, display={First pair: yes, fully equivalent. Combining gives exactly $\frac{2x}{x^2-4}$, and both forms accept every input except $x=2,-2$, so they agree everywhere either one is defined. Second pair: \emph{not} interchangeable. They agree at every input both accept, but $\frac{3}{x+3}$ also accepts $x=3$ (giving $\frac12$) while the difference does not --- so as functions one has a hole at $\left(3,\frac12\right)$ and the other does not.}]{\linewidth}{\gradedlinesheight}}
\end{enumerate}
\end{notesbox}

\begin{extensionbox}
\textbf{Part 1 --- Three students, three answers.} Each student combined $\dfrac{5}{x-1}-\dfrac{2}{x+1}$
and got something different. Decide who is right, and name each error.
\begin{center}
\renewcommand{\arraystretch}{1.4}
\begin{tabularx}{\linewidth}{@{}l l X@{}}
\hline
\textbf{Student} & \textbf{Answer} & \textbf{Right or wrong --- and why?} \\
\hline
Gia   & $\dfrac{3}{2x}$ & \blank{5.4cm} \\
Hana  & $\dfrac{3}{x-1}$, $x\neq1,-1$ & \blank{5.4cm} \\
Ivo   & $\dfrac{3x+7}{x^2-1}$, no restrictions & \blank{5.4cm} \\
\hline
\end{tabularx}
\end{center}
\vspace{2pt}
Hana's answer is the dangerous one: her mistake produced a numerator that \emph{cancels beautifully}.
Show her work, find the exact step that went wrong, and then test her answer and the correct one at
$x=0$. \writelines{2}

\vspace{4pt}
\textbf{Part 2 --- Design one.} Write a \emph{sum} of two rational expressions whose simplified value is
$\dfrac{5}{x-1}$ and whose restrictions are exactly $x\neq1$ and $x\neq-2$.
\begin{center}
\blank{3.6cm} $+$ \blank{3.6cm}
\end{center}
\begin{enumerate}[label=(\alph*), leftmargin=1.9em, itemsep=5pt]
  \item Which of your denominators had to be a quadratic, and why could they not both be linear?
  \item Which excluded value is a hole on the graph of your expression, and which is a wall?
\end{enumerate}
\writelines{3}
\end{extensionbox}

\begin{spiralbox}
\textbf{Coming next --- Lesson 4.4: The Rational Parent Function \& Transformations.} For four lessons
the rational functions you have met have been \emph{expressions} to simplify. Now they go back on the
grid. Every one of them is a relative of a single parent, $y=\dfrac{1}{x}$ --- the hyperbola with a wall
at $x=0$ and a floor at $y=0$ --- and the same four transformations you used on lines, parabolas, and
absolute values ($f(x)+k$, $f(x+k)$, $kf(x)$, $f(kx)$) apply to it. The new twist: shifting a rational
function \emph{moves its asymptotes}, so the walls and floors you learned to read in 4.0 stop sitting on
the axes.
\end{spiralbox}

\end{document}
